\section{Mathematical derivations and diagnostic identities}
\label{app:theory}
\subsection{Conditional effective-matrix moments}
Let $A,B,P$ be fixed conditional states after the deterministic signal part
of a round. Assume independent entries of $Z_A\in\R^{M\times r}$ and
$Z_B\in\R^{r\times d}$ have mean zero and variance $\tau^2$, with the
two matrices independent. In the implemented mechanism they are Gaussian.
Write $D_1=AZ_B$, $D_2=Z_AB$, and $D_3=Z_AZ_B$.

Using $\E[Z_BZ_B^T]=d\tau^2I_r$ gives
\begin{equation}
 \E\norm{D_1}_F^2=d\tau^2\norm{A}_F^2,
 \qquad \E\norm{D_2}_F^2=M\tau^2\norm{B}_F^2.
\end{equation}
Every entry of $Z_AZ_B$ is a sum of $r$ products of independent centered
variables, with variance $r\tau^4$. Summing over its $Md$ entries gives
$\E\norm{D_3}_F^2=Mdr\tau^4$. Each cross expectation vanishes by
conditioning on one centered noise matrix. For example, conditional on
$Z_B$, the inner product $\ip{AZ_B}{Z_AZ_B}_F$ is linear in $Z_A$ and
has expectation zero. Expanding the squared norm proves
\cref{eq:qmoment}; mutual independence of all three products is unnecessary.

\subsection{Score moments and pairwise margins}
For the first score term, $PD_1^T=PZ_B^TA^T$, the identity
$\E[Z_B^TA^TAZ_B]=\tau^2\norm{A}_F^2I_d$ yields
\begin{equation}
 \E\norm{PD_1^T}_F^2=\tau^2\norm{A}_F^2\norm{P}_F^2.
\end{equation}
For $PD_2^T=PB^TZ_A^T$ the expectation is
$M\tau^2\norm{PB^T}_F^2$. Conditioning successively on both noises
gives $Mr\tau^4\norm{P}_F^2$ for $PD_3^T$. Cross expectations again
vanish, giving \cref{eq:smoment}.

For distinct items $i,j$ and fixed user $p$, the noisy difference between
their scores has a centered shock with variance
\begin{equation}
 V_{ij}=2\tau^2\norm{Bp}^2+
 \tau^2\norm{A_i-A_j}^2\norm{p}^2+
 2r\tau^4\norm{p}^2.
 \label{eq:marginvar}
\end{equation}
The factor two arises because the two independent rows of $Z_A$ are
subtracted. Shared $Z_B$ affects both items through their factor difference.
The full shock is not generally Gaussian. For clean nonzero margin $m$,
Chebyshev's inequality supplies the weak bound
\begin{equation}
 \Pr(\text{sign flip})\leq
 \Pr(|\text{shock}|\geq |m|)\leq \min(1,V_{ij}/m^2).
\end{equation}
A variance alone does not justify a Gaussian tail formula for Two.

For FixedB the shock is $(Z_{A,i}-Z_{A,j})Bp$, exactly Gaussian with
variance $2\tau^2\norm{Bp}^2$. Its fixed-margin sign-flip probability is
$\Phi(-|m|/(\sqrt{2}\tau\norm{Bp}))$ when the denominator is positive.
This formula conditions on the current factors and user state; it is not
an expression for full-catalog NDCG or for accumulated training loss.

\subsection{Scale changes, balancing, and clipping}
The transformation $(A,B)\mapsto(cA,B/c)$ preserves $Q$ for positive $c$.
The linear portion of expected matrix-noise energy becomes
\begin{equation}
 \tau^2\left(\frac{M\norm{B}_F^2}{c^2}
                       +dc^2\norm{A}_F^2\right).
\end{equation}
Differentiating gives the minimizer
$c^4=M\norm{B}_F^2/(d\norm{A}_F^2)$ when both factors are nonzero.
For scores, the corresponding minimizer satisfies
$c^4=M\norm{PB^T}_F^2/(\norm{A}_F^2\norm{P}_F^2)$.
An equal-factor-norm convention need not minimize either rectangular or
user-weighted expression. The bilinear $\tau^4$ term is unchanged by
this scalar gauge transformation.

These conditional optima are not a training prescription. Changing factor
scale also changes client gradients and which updates are clipped. A rule
chosen from private $P$ would require a privacy treatment of that information.
Balancing after the noisy update preserves its already-noisy $AB$ and
cannot remove the realized effective shock. The formulas explain an
interpretation problem, not a free improvement of the privacy mechanism.

Separate factor averaging has an additional distinction. With client means
$\bar A,\bar B$,
\begin{equation}
 \frac1K\sum_{u=1}^K A_uB_u
 =\bar A\bar B+\frac1K\sum_{u=1}^K(A_u-\bar A)(B_u-\bar B).
\end{equation}
Thus averaging factors and averaging products generally differ by a
cross-client covariance term. FixedB removes this discrepancy, but E1
does not independently intervene on it. It cannot be isolated as the
reason for a utility difference from the existing comparisons.

\subsection{FixedB score capacity and regularization}
If $BB^T=I_r$, define $v_u=Bp_u$. The scores are $A_iv_u$, with a score
matrix of rank at most $r$. Conversely any desired $v_u\in\R^r$ has
preimage $p_u=B^Tv_u$, and every other preimage adds a vector in the null
space of $B$. Orthogonality implies
\begin{equation}
 \min_{p:Bp=v}\norm{p}^2=\norm{v}^2,
 \qquad \norm{A_iB}^2=\norm{A_i}^2.
\end{equation}
Projecting the local gradient through $B$ therefore gives ordinary rank-$r$
BPR dynamics for the score-relevant vector, under the same effective-row
penalty, up to floating-point behavior. A null-space component does not
change scores. FixedB does not encode public item semantics; it defines
a public coordinate map for private collaborative factors.

\subsection{A product-preserving coupling null example}
The pairs $(A,B)$ and $(-A,-B)$ represent exactly the same matrix. Applying
identical raw-coordinate Gaussian draws to both gives noisy products whose
difference is
\begin{equation}
 (A+Z_A)(B+Z_B)-(-A+Z_A)(-B+Z_B)
 =2(AZ_B+Z_AB).
\end{equation}
The expected squared distance is
$4\tau^2(d\norm{A}_F^2+M\norm{B}_F^2)$, although the two noisy products
have the same marginal distribution. Coupling the second draw as
$(-Z_A,-Z_B)$ instead makes the products identical. A retained 10,000-draw
synthetic check gives measured energy $6.341142$ against theoretical
$6.345313$, with Monte Carlo SE $0.012443$; aligned distance is exactly zero.
This elementary example makes the coupling issue explicit without claiming
a new theorem or an accuracy gain.

\subsection{Exact conditional margin law and numerical evaluation}
Let $\tilde X=X+a$, $\tilde Y=Y+b$ in \cref{eq:tailmargin}. Then
$(\tilde X+\tilde Y)/\sqrt2$ and $(\tilde X-\tilde Y)/\sqrt2$ are independent
normal vectors with identity covariance and means $(a+b)/\sqrt2$ and
$(a-b)/\sqrt2$. The identity
$\tilde X^T\tilde Y=(\norm{\tilde X+\tilde Y}^2-
\norm{\tilde X-\tilde Y}^2)/4$ proves \cref{eq:chisquare}. The first
three cumulants are
\begin{equation}
 \E L=m,\qquad \Var(L)=V_{\mathrm{lin}}+rk,
 \qquad \kappa_3(L)=6km.
\end{equation}
For positive $m$, the sign probability can be evaluated as
$\int_0^\infty F_U(v)f_V(v)\,dv$. The implementation controls omitted
tails, uses a square-root variable near zero, and compares against direct
Monte Carlo in synthetic cases. Positive rescaling of $p$ multiplies both
margin and shock by the same positive amount, leaving sign-flip probability
unchanged. The real-state calculation normalizes this scale for numerical
stability; it does not change the ranking question.

\subsection{Analytic Gaussian calibration for popularity}
Under add/remove adjacency, \cref{eq:pop} changes the unnoised sum by a
vector of norm at most $\sqrt{20}$. For noise SD equal to sensitivity times
$\sigma_{\mathrm{pop}}$, the analytic Gaussian condition is
\begin{equation}
 \delta=\Phi\left(\frac{1}{2\sigma_{\mathrm{pop}}}
                           -\epsilon\sigma_{\mathrm{pop}}\right)
 -e^\epsilon\Phi\left(-\frac{1}{2\sigma_{\mathrm{pop}}}
                           -\epsilon\sigma_{\mathrm{pop}}\right).
\end{equation}
Solving numerically yields multipliers approximately $0.600229$, $1.081162$,
$1.993812$, and $3.730632$ for targets 8, 4, 2, and 1. This is one release
with $q=1$, not 50 compositions of a subsampled update. Releasing all
noise-seed versions would require additional accounting. The derivation
uses the established analytic Gaussian mechanism \citep{balle2018gaussian}.
